\section*{Chapitre \ref{ch:b3:rings} --- Anneaux et arithmétique}

\begin{solution}{exo:b3:rings:1}
(a) L'évaluation $f \colon \Z[X] \to \Z[\iu]$, $P \mapsto
P(\iu)$, est un morphisme d'anneaux surjectif ($a + bX \mapsto a
+ b\iu$). Noyau~: diviser $P$ par le polynôme \emph{unitaire}
$X^2 + 1$ dans $\Z[X]$~: $P = (X^2 + 1)Q + (bX + a)$ avec $a, b
\in \Z$~; alors $P(\iu) = a + b\iu = 0$ ssi $a = b = 0$. Donc
$\ker f = (X^2+1)$ et le~\cref{thm:b3:rings:firstiso} conclut.

(b) Le même calcul à coefficients dans $\R$~: $\R[X]/(X^2+1)
\cong \C$ --- c'est la \emph{construction} la plus propre de
$\C$.

(c) $X^2 + X + 1$ n'a pas de racine dans $\mathbb F_2$ ($0, 1
\mapsto 1$), donc, de degré $2$, il est \omterm{def:b3:rings:divisibility}{irréductible}~: le
quotient $\mathbb F_4 = \mathbb F_2[X]/(X^2+X+1)$ est un corps
(\cref{prop:b3:rings:primemaximal}~; $(P)$ \omterm{def:b3:rings:primemaximal}{maximal} dans $K[X]$
quand $P$ est \omterm{def:b3:rings:divisibility}{irréductible}, car $K[X]$ est un \omterm{def:b3:rings:pidufd}{PID}~: un \omterm{def:b3:rings:ideal}{idéal}
$(D) \supseteq (P)$ signifie $D \mid P$). Ses quatre éléments
sont $0, 1, \omega, \omega + 1$ où $\omega = \bar X$, avec
$\omega^2 = \omega + 1$. Table de multiplication (éléments non
nuls)~:
\[
\omega \cdot \omega = \omega + 1, \qquad
\omega(\omega + 1) = \omega^2 + \omega = 1, \qquad
(\omega+1)^2 = \omega^2 + 1 = \omega .
\]
Les éléments non nuls forment un groupe cyclique d'ordre $3$
engendré par $\omega$.
\end{solution}

\begin{solution}{exo:b3:rings:2}
(a) Soit $p$ \omterm{def:b3:rings:divisibility}{irréductible} dans un \omterm{def:b3:rings:pidufd}{UFD} et $p \mid ab$, disons $ab
= pc$, avec $a, b \neq 0$ (sinon trivial). Si $a$ ou $b$ est une
unité, $p$ divise l'autre. Sinon factoriser $a$, $b$ et $c$ en
\omterm{def:b3:rings:divisibility}{irréductibles}~: les deux factorisations de $ab$,
\[
(\text{facteurs de } a)(\text{facteurs de } b) = p \cdot
(\text{facteurs de } c),
\]
doivent s'accorder à l'ordre et aux \omterm{def:b3:rings:divisibility}{associés} près~: $p$ est
associé d'un certain facteur \omterm{def:b3:rings:divisibility}{irréductible} de $a$ ou de $b$, donc
le divise.

(b) Soit $A$ un anneau intègre fini et $x \neq 0$. L'application
$y \mapsto xy$ est injective ($xy = xy' \Rightarrow x(y - y') =
0 \Rightarrow y = y'$), donc surjective ($A$ fini)~: $1 = xy$
pour un certain $y$.

(c) Si $\mathfrak p$ est premier dans un anneau fini $A$, alors
$A/\mathfrak p$ est un anneau intègre fini, donc un corps par
(b), d'où $\mathfrak p$ est \omterm{def:b3:rings:primemaximal}{maximal}
(\cref{prop:b3:rings:primemaximal}).
\end{solution}

\begin{solution}{exo:b3:rings:3}
Normes~: $N(3) = 9$, $N(1 \pm \iu\sqrt5) = 6$, $N(2 \pm
\iu\sqrt5) = 9$. Les équations $x^2 + 5y^2 = 2$ et $x^2 + 5y^2 =
3$ n'ont pas de solutions entières, donc aucun élément n'a norme
$2$ ou $3$. Une factorisation propre de $3$ demanderait deux
facteurs de norme $3$~: impossible --- $3$ est \omterm{def:b3:rings:divisibility}{irréductible}. Une
factorisation propre de $1 \pm \iu\sqrt5$ (norme $6$) demanderait
des facteurs de normes $2, 3$~: impossible. De même pour $2 \pm
\iu\sqrt5$ (norme $9$~: les facteurs auraient norme $3$).
Maintenant
\[
9 = 3 \cdot 3 = (2 + \iu\sqrt5)(2 - \iu\sqrt5),
\]
deux factorisations en \omterm{def:b3:rings:divisibility}{irréductibles}. Elles sont véritablement
différentes~: les unités sont $\pm 1$ (norme $1$), et $2 \pm
\iu\sqrt5 \neq \pm 3$. Donc l'unicité échoue --- tandis que
l'existence des factorisations vaut dans $\Z[\iu\sqrt5]$
(\cref{exo:b3:rings:10}(c))~: la non-factorialité ici est purement
un échec d'unicité. (De façon cohérente,
\cref{prop:b3:rings:primeirred}~: ces \omterm{def:b3:rings:divisibility}{irréductibles} ne sont pas
premiers.)
\end{solution}

\begin{solution}{exo:b3:rings:4}
(a) Soient $a, b \in \Z[\iu]$, $b \neq 0$, et $a/b = x + \iu y
\in \C$. Choisir des entiers $m, n$ avec $\abs{x - m} \leq
\frac12$, $\abs{y - n} \leq \frac12$, et poser $q = m + \iu n$,
$r = a - bq$. Alors
\[
N(r) = N(b)\,\abs*{\tfrac ab - q}^2
\leq N(b)\Bigl(\tfrac14 + \tfrac14\Bigr) = \tfrac{N(b)}2 < N(b).
\]
Donc $N$ est une fonction euclidienne ($N(r) < N(b)$ ou $r =
0$).

(b) Si $uv = 1$ alors $N(u)N(v) = 1$ avec $N(u) \in \N$~: $N(u)
= 1$, i.e.\ $x^2 + y^2 = 1$~: $u \in \{\pm 1, \pm\iu\}$~;
réciproquement ce sont des unités.

(c) Pour $\Z[\iu\sqrt2]$~: le même arrondi donne $\abs{a/b -
q}^2 \leq \frac14 + \frac{2}4 = \frac34 < 1$~: \omterm{def:b3:rings:pidufd}{euclidien}~;
unités~: $x^2 + 2y^2 = 1$ donne $\pm 1$. Pour
$\Z[\iu\sqrt5]$ la borne devient $\frac14 + \frac54 = \frac32 >
1$~: l'argument d'arrondi échoue --- et doit échouer, car
$\Z[\iu\sqrt5]$ n'est même pas un \omterm{def:b3:rings:pidufd}{UFD}
(\cref{exo:b3:rings:3}), tandis qu'\omterm{def:b3:rings:pidufd}{euclidien} impliquerait \omterm{def:b3:rings:pidufd}{UFD}
(\cref{thm:b3:rings:euclideanpid,thm:b3:rings:pidufd}).
\end{solution}

\begin{solution}{exo:b3:rings:5}
(a) Si $x^n = 0$~:
\[
(1 + x)\bigl(1 - x + x^2 - \dots + (-1)^{n-1}x^{n-1}\bigr)
= 1 + (-1)^{n-1}x^n = 1 .
\]
(b) Dans un anneau intègre, $\deg(PQ) = \deg P + \deg Q$~; $PQ =
1$ force $\deg P = \deg Q = 0$ et $P, Q \in A^\times$~:
$A[X]^\times = A^\times$. Sur $\Z/4\Z$~: $(1 + 2X)^2 = 1 + 4X +
4X^2 = 1$, donc $1 + 2X$ est une unité de degré $1$ (ici $2$ est
nilpotent~; comparer (a)).

(c) $e^2 = e$ donne $e(e - 1) = 0$, donc $e \in \{0, 1\}$ dans un
anneau intègre. Si $x^n = 0$ avec $n \geq 1$ minimal et $x \ne
0$, alors $n \geq 2$ et $x \cdot x^{n-1} = 0$ avec les deux
facteurs non nuls~: contradiction.
\end{solution}

\begin{solution}{exo:b3:rings:6}
(a) $K[X,Y]/(X,Y) \cong K$ (évaluer en $(0,0)$)~: un corps, donc
$(X,Y)$ est \omterm{def:b3:rings:primemaximal}{maximal}. Si $(X, Y) = (P)$~: $P \mid X$ force (degrés
en $Y$) $P \in K[X]$, et $P \mid Y$ force alors $P \in K$~; $P =
0$ est absurde et $P \in K^\times$ donnerait $(P) = K[X,Y]$,
contredisant le caractère propre ($K[X,Y]/(X,Y) \cong K \neq
0$). Donc $(X,Y)$ n'est pas principal.

(b) L'évaluation $P(X, Y) \mapsto P(X, X^2)$ envoie $K[X,Y]$ sur
$K[X]$~; son noyau est $(Y - X^2)$~: en divisant par le polynôme
unitaire en $Y$ $Y - X^2$, $P = (Y - X^2)Q + R(X)$, et $P(X,
X^2) = R(X)$. Donc $K[X,Y]/(Y - X^2) \cong K[X]$ --- l'anneau
des coordonnées d'une parabole, isomorphe à celui d'une droite.

L'évaluation $P(X, Y) \mapsto P(X, X^{-1})$ envoie $K[X, Y]$ sur
l'anneau $K[X, X^{-1}]$ des polynômes de Laurent. Son noyau
contient $(XY - 1)$~; réciproquement, modulo $XY - 1$ toute
classe a un représentant $R = \sum_{n \geq 0} a_nX^n +
\sum_{m \geq 1} b_m Y^m$ (remplacer chaque produit $XY$ par $1$
répétitivement), et $R(X, X^{-1}) = \sum a_n X^n + \sum b_m
X^{-m} = 0$ force tous les $a_n = b_m = 0$. D'où $K[X, Y]/(XY -
1) \cong K[X, X^{-1}]$ --- l'anneau des coordonnées d'une
hyperbole~: la droite privée d'un point.

(c) $(Y - X^2)$ est premier (le quotient $K[X]$ est intègre) mais
non \omterm{def:b3:rings:primemaximal}{maximal} ($K[X]$ n'est pas un corps~; concrètement $(Y - X^2)
\subsetneq (Y - X^2,\, X) \subsetneq K[X,Y]$).
\end{solution}

\begin{solution}{exo:b3:rings:7}
\emph{$X^5 - 12X^3 + 36X - 12$}~: \omterm{thm:b3:rings:criteria}{Eisenstein} en $p = 3$ ($3 \mid
12, 36, 12$~; $9 \nmid 12$~; $3 \nmid 1$)~: \omterm{def:b3:rings:divisibility}{irréductible}. (En $p
= 2$ \omterm{thm:b3:rings:criteria}{Eisenstein} échoue~: $4 \mid 12$.)

\emph{$X^4 + X + 1$}~: réduire mod $2$. Pas de racine dans
$\mathbb F_2$~; le seul quadratique \omterm{def:b3:rings:divisibility}{irréductible} sur $\mathbb
F_2$ est $X^2 + X + 1$, et $(X^2+X+1)^2 = X^4 + X^2 + 1 \neq X^4
+ X + 1$. Donc $X^4 + X + 1$ est \omterm{def:b3:rings:divisibility}{irréductible} sur $\mathbb F_2$,
d'où sur $\Q$ (\cref{thm:b3:rings:criteria}(1)~; il est
unitaire).

\emph{$X^4 + 4$}~: réductible --- l'identité de Sophie Germain,
$X^4 + 4 = (X^2 - 2X + 2)(X^2 + 2X + 2)$.

\emph{$X^4 + 1$}~: décaler, $(X+1)^4 + 1 = X^4 + 4X^3 + 6X^2 +
4X + 2$~: \omterm{thm:b3:rings:criteria}{Eisenstein} en $2$. Une factorisation de $X^4+1$ se
décalerait en une de $(X+1)^4 + 1$~: \omterm{def:b3:rings:divisibility}{irréductible}.

\emph{$X^3 - X - 1$}~: un cubique est réductible sur $\Q$ ssi il
a une racine rationnelle~; une racine rationnelle d'un polynôme
entier unitaire est un entier divisant le terme constant
(théorème des racines rationnelles~: si $(p/q)$ en termes
réduits est une racine, $q \mid 1$, $p \mid -1$), et $\pm 1$ ne
sont pas racines ($-1$ et $-1$)~: \omterm{def:b3:rings:divisibility}{irréductible}.
\end{solution}

\begin{solution}{exo:b3:rings:8}
(a) Pour $\gcd(m, n) = 1$, le~\cref{thm:b3:rings:crt} donne un
isomorphisme d'anneaux $\Z/mn\Z \cong \Z/m\Z \times \Z/n\Z$. Un
élément d'un anneau produit est une unité ssi les deux
coordonnées le sont, donc $(\Z/mn\Z)^\times \cong
(\Z/m\Z)^\times \times (\Z/n\Z)^\times$ et $\varphi(mn) =
\varphi(m)\varphi(n)$. Pour une puissance de premier, les non
unités de $\Z/p^k\Z$ sont les classes des multiples de $p$~:
$\varphi(p^k) = p^k - p^{k-1}$. D'où
\[
\varphi(n) = \prod_i \bigl(p_i^{\alpha_i} -
p_i^{\alpha_i-1}\bigr) = n \prod_{p \mid n}\Bigl(1 -
\frac1p\Bigr).
\]

(b) $M = 7 \cdot 11 \cdot 13 = 1001$. Idempotents~: $e_1 \equiv
(1, 0, 0)$~: $143 = 11\cdot13 \equiv 3 \pmod 7$ et $3 \cdot 5
\equiv 1$~: $e_1 = 143 \cdot 5 = 715$. $e_2$~: $91 \equiv 3
\pmod{11}$, $3 \cdot 4 \equiv 1$~: $e_2 = 91\cdot4 = 364$.
$e_3$~: $77 \equiv -1 \pmod{13}$~: $e_3 = 77 \cdot 12 = 924$.
Alors
\[
x \equiv 2\,e_1 + 5\,e_2 + 1\,e_3 = 1430 + 1820 + 924 = 4174
\equiv 170 \pmod{1001},
\]
et en effet $170 = 24\cdot7 + 2 = 15\cdot11 + 5 = 13\cdot13 +
1$.
\end{solution}

\begin{solution}{exo:b3:rings:9}
(a) Si $x^n = 0$ et $y^m = 0$, le développement binomial de
$(x+y)^{n+m}$ a tout terme $x^iy^j$ avec $i + j = n + m$, donc
$i \geq n$ ou $j \geq m$~: chaque terme s'annule, et $x + y$ est
nilpotent~; $(ax)^n = a^nx^n = 0$~: $\operatorname{Nil}(A)$ est
un \omterm{def:b3:rings:ideal}{idéal}. Si $\mathfrak p$ est premier et $x^n = 0 \in
\mathfrak p$, une récurrence sur $n$ donne $x \in \mathfrak p$
($x \cdot x^{n-1} \in \mathfrak p$).

(b) Soit $a \notin \operatorname{Nil}(A)$ et $S = \{a^n : n \geq
1\}$, donc $0 \notin S$. L'ensemble $\mathcal E$ des idéaux
disjoints de $S$ contient $(0)$ et est inductif (l'union d'une
chaîne d'idéaux disjoints de $S$ est un \omterm{def:b3:rings:ideal}{idéal} disjoint de $S$)~:
Zorn fournit $\mathfrak p \in \mathcal E$ \omterm{def:b3:rings:primemaximal}{maximal}. $\mathfrak
p$ est propre ($a \notin \mathfrak p$, car $a \in S$).
Primalité~: soient $x, y \notin \mathfrak p$. Par maximalité,
$\mathfrak p + (x)$ et $\mathfrak p + (y)$ rencontrent $S$~:
$a^m \in \mathfrak p + (x)$, $a^n \in \mathfrak p + (y)$. En
multipliant, $a^{m+n} \in \mathfrak p + (xy)$. Si $xy \in
\mathfrak p$, alors $a^{m+n} \in \mathfrak p \cap S$~:
absurde. Donc $xy \notin \mathfrak p$ --- contraposée de la
primalité. D'où tout élément non nilpotent évite un certain
\omterm{def:b3:rings:primemaximal}{idéal premier}~; avec (a), $\operatorname{Nil}(A) =
\bigcap_{\mathfrak p} \mathfrak p$.
\end{solution}

\begin{solution}{exo:b3:rings:10}
(a) La chaîne $\ker f \subseteq \ker f^2 \subseteq \cdots$ se
stabilise (\cref{prop:b3:rings:noethacc})~: $\ker f^n = \ker
f^{n+1}$ pour un certain $n$. Soit $x \in \ker f$. Comme $f$,
donc $f^n$, est surjective, $x = f^n(y)$ pour un certain $y$~;
alors $f^{n+1}(y) = f(x) = 0$, donc $y \in \ker f^{n+1} = \ker
f^n$, i.e.\ $x = f^n(y) = 0$.

(b) $I_n = \{f \in \mathcal C(\intcc01, \R) : f\restriction_{
\intcc0{1/n}} = 0\}$ est un \omterm{def:b3:rings:ideal}{idéal}, et $I_n \subseteq I_{n+1}$.
L'inclusion est stricte~: $x \mapsto \max\bigl(0, x -
\frac1{n+1}\bigr)$ s'annule sur $\intcc0{\frac1{n+1}}$ mais pas
sur $\intcc0{\frac1n}$. Une chaîne strictement croissante
infinie contredit la~\cref{prop:b3:rings:noethacc}.

(c) Supposer que l'ensemble des non nuls non inversibles
n'admettant pas de factorisation en \omterm{def:b3:rings:divisibility}{irréductibles} est non vide.
La famille correspondante d'idéaux $\{(a)\}$ a un élément
\omterm{def:b3:rings:primemaximal}{maximal} $(a)$ (\cref{prop:b3:rings:noethacc}). L'élément $a$
n'est pas \omterm{def:b3:rings:divisibility}{irréductible} (un \omterm{def:b3:rings:divisibility}{irréductible} est sa propre
factorisation), donc $a = bc$ avec $b, c$ non inversibles~; $(a)
\subseteq (b)$ est strict (car $(a) = (b)$ donnerait $b = ad$, $a
= adc$, donc $dc = 1$~: $c$ une unité), de même $(a) \subsetneq
(c)$. Par maximalité, $b$ et $c$ se factorisent tous deux en
\omterm{def:b3:rings:divisibility}{irréductibles}~; en concaténant les facteurs de $a$~:
contradiction. Appliqué à $\Z[\iu\sqrt5]$ --- \omterm{def:b3:rings:noetherian}{noethérien} comme
quotient de $\Z[X]$ (\cref{thm:b3:rings:hilbert},
$\Z[\iu\sqrt5] \cong \Z[X]/(X^2+5)$) --- cela montre que les
factorisations existent là~; \cref{exo:b3:rings:3} montrait que
c'est l'unicité qui échoue.
\end{solution}

\begin{solution}{exo:b3:rings:11}
(a) \cref{exo:b3:rings:7}~: décalage et \omterm{thm:b3:rings:criteria}{Eisenstein} en $2$.

(b) Mod $2$~: $X^4 + 1 = (X + 1)^4$. Soit maintenant $p$ impair.
Les carrés forment un sous-groupe d'indice $2$ dans
$(\Z/p\Z)^\times$~: le morphisme $x \mapsto x^2$ a pour noyau
$\{\pm 1\}$ (deux éléments~: $X^2 - 1$ a au plus $2$ racines
dans un corps, et $1 \neq -1$ pour $p$ impair), donc son image a
$\frac{p-1}2$ éléments. Par conséquent, le produit de deux non
carrés est un carré (dans le \omterm{thm:b3:groups:quotient}{groupe quotient} d'ordre $2$,
$\overline{xy} = \bar x\bar y$). D'où \emph{au moins l'un de
$-1$, $2$, $-2$ est un carré mod $p$} (si $-1$ et $2$ ne le
sont pas, $-2 = (-1)\cdot 2$ l'est). Dans chaque cas $X^4 + 1$
se factorise mod $p$~:
\begin{itemize}
  \item $-1 = c^2$~: \quad $X^4 + 1 = X^4 - c^2 = (X^2 -
        c)(X^2 + c)$~;
  \item $2 = c^2$~: \quad $(X^2 + cX + 1)(X^2 - cX + 1) = X^4 +
        (2 - c^2)X^2 + 1 = X^4 + 1$~;
  \item $-2 = c^2$~: \quad $(X^2 + cX - 1)(X^2 - cX - 1) = X^4
        - (c^2 + 2)X^2 + 1 = X^4 + 1$.
\end{itemize}
Donc $X^4+1$ est réductible modulo tout premier, pourtant
\omterm{def:b3:rings:divisibility}{irréductible} sur $\Q$~: le critère de réduction
(\cref{thm:b3:rings:criteria}(1)) détecte l'irréductibilité
mais son échec ne prouve rien.

(Pour la raison structurelle~: $p^2 - 1 = (p-1)(p+1)$ est un
produit de deux nombres pairs consécutifs, donc $8 \mid p^2 -
1$~; le groupe cyclique $\mathbb F_{p^2}^\times$ (cyclicité
démontrée dans le~\cref{ch:b3:galois}) contient alors un élément
$\zeta$ d'ordre $8$, racine de $X^4 + 1$~; son polynôme minimal
sur $\mathbb F_p$ divise $X^4+1$ et a degré $\leq 2$ ---
$X^4+1$ ne peut jamais être \omterm{def:b3:rings:divisibility}{irréductible} mod $p$.)
\end{solution}

\begin{solution}{exo:b3:rings:12}
(a) $(1-e)^2 = 1 - 2e + e^2 = 1 - e$. L'application
$\varphi(x) = (ex, (1-e)x)$ est additive et multiplicative vers
le produit des deux idéaux~: $exey = e^2xy = e(xy)$, et $Ae$ est
un anneau commutatif d'unité $e$ ($e\cdot ex = ex$). Injectif~:
$ex = 0$ et $(1-e)x = 0$ s'additionnent en $x = 0$. Surjectif~:
$(ea, (1-e)b)$ est l'image de $ea + (1-e)b$ (calculer les deux
composantes en utilisant $e(1-e) = 0$). Les unités s'envoient
correctement en paires de style $(1, 0)$~: $\varphi(1) = (e,
1-e)$, l'unité du produit.

(b) Dans un anneau intègre, $e(e - 1) = 0$ force $e \in \{0,
1\}$~: seulement des idempotents triviaux. Dans $\Z/12\Z$,
résoudre $e^2 \equiv e$~: $e \in \{0, 1, 4, 9\}$. La paire non
triviale $\{4, 9\}$~: $4 + 9 = 13 \equiv 1$, $4\cdot9 = 36
\equiv 0$, et $\Z/12\Z\cdot4 = \{0, 4, 8\} \cong \Z/3\Z$ (unité
$4$), $\Z/12\Z\cdot9 = \{0, 3, 6, 9\} \cong \Z/4\Z$ (unité
$9$)~: le scindage CRT $\Z/12\Z \cong \Z/4\Z\times\Z/3\Z$, avec
$9 \leftrightarrow (1, 0)$ et $4 \leftrightarrow (0, 1)$.

(c) Sous l'isomorphisme CRT $\Z/n\Z \cong \prod_i
\Z/p_i^{a_i}\Z$, l'élément $e_i$ aux congruences énoncées
correspond au $n$-uplet avec $1$ dans la case $i$ et $0$ ailleurs~:
les idempotents élémentaires du produit. Réciproquement une
famille complète d'idempotents orthogonaux ($e_ie_j = 0$ pour $i
\neq j$, $\sum e_i = 1$) réassemble la décomposition en produit
par (a), par récurrence. Les idempotents sont aux anneaux ce que
les projections orthogonales sont aux espaces de Hilbert
(\cref{ch:b3:hilbert})~: les coordonnées d'une décomposition
directe interne.
\end{solution}

\begin{solution}{pb:b3:rings:1}
\textbf{1.} $N(z) = z\bar z$, donc $N(zw) = zw\overline{zw} =
z\bar z\, w \bar w = N(z)N(w)$. Si $uv = 1$~: $N(u)N(v) = 1$ dans
$\N$, donc $N(u) = 1$, i.e.\ $u \in \{\pm 1, \pm \iu\}$~; les
quatre sont des unités. Brahmagupta~: $(a^2+b^2)(c^2+d^2) =
N\bigl((a + \iu b)(c + \iu d)\bigr) = (ac - bd)^2 + (ad +
bc)^2$.

\textbf{2.} Si $z = ab$, alors $N(z) = N(a)N(b)$ est premier,
donc $N(a) = 1$ ou $N(b) = 1$~: un facteur est une unité. Comme
$N(z) > 1$, $z$ n'est ni zéro ni une unité~: \omterm{def:b3:rings:divisibility}{irréductible} --- et
premier, car $\Z[\iu]$ est un \omterm{def:b3:rings:pidufd}{UFD}
(\cref{thm:b3:rings:euclideanpid,thm:b3:rings:pidufd,%
lem:b3:rings:bezout}).

\textbf{3.} $\pi$ divise $N(\pi) = \pi\bar\pi \geq 2$, un
entier~; en factorisant $N(\pi)$ en nombres premiers et en
utilisant que $\pi$ est premier, $\pi \mid p$ pour un certain
nombre premier $p$. Si aussi $\pi \mid q \neq p$~: Bézout dans
$\Z$ donne $1 = up + vq$, donc $\pi \mid 1$ --- absurde~: $p$
est unique. De $p = \pi\gamma$~: $p^2 = N(p) = N(\pi)N(\gamma)$
avec $N(\pi) \neq 1$, donc $N(\pi) \in \{p, p^2\}$.

\textbf{4.} Soit $\pi$ un premier de Gauss divisant $p$, $p =
\pi\gamma$. Si $N(\pi) = p^2$~: $N(\gamma) = 1$, donc $p$ est
associé de $\pi$, lui-même un premier de Gauss~; et aucun
premier de Gauss n'a norme $p$ (si $N(\rho) = p$ alors $\rho
\mid \rho\bar\rho = p$, et $p$ premier dans $\Z[\iu]$ forcerait
$\rho$ associé à $p$, donnant $N(\rho) = N(p) = p^2 \neq p$).
Si $N(\pi) = p$~: en écrivant $\pi = a + \iu b$, $p =
\pi\bar\pi = a^2 + b^2$.

\textbf{5.} Dans le groupe abélien $(\Z/p\Z)^\times$, apparier
chaque élément avec son inverse. Les éléments auto-inverses sont
les racines de $X^2 - 1$~: exactement $\pm 1$ (au plus deux
racines dans un corps). Le produit de tous les éléments est alors
$1 \cdot (-1) \cdot \prod (\text{paires } k k^{-1}) = -1$~:
$(p-1)! \equiv -1 \pmod p$. (Pour $p = 2$~: $1! \equiv -1$.)

\textbf{6.} Écrire $(p-1)! = \prod_{k=1}^m k \cdot
\prod_{k=m+1}^{p-1}k$ avec $m = \frac{p-1}2$. Dans le second
produit substituer $k = p - j$, $j = 1, \dots, m$~: modulo $p$,
$\prod_{j=1}^m (p - j) \equiv (-1)^m m!$. D'où $-1 \equiv
(-1)^m (m!)^2$, i.e.\ $(m!)^2 \equiv (-1)^{m+1} \pmod p$.

\textbf{7.} Si $p \equiv 1 \pmod 4$, $m$ est pair et la question
6 donne $(m!)^2 \equiv -1$~: une racine carrée de $-1$.
Réciproquement, si $x^2 \equiv -1 \pmod p$ ($p$ impair), alors
$x^4 = 1 \neq x^2$~: $x$ a ordre $4$ dans $(\Z/p\Z)^\times$,
donc $4 \mid p - 1$ (Lagrange). Et $p = 2$~: $1^2 = 1 \equiv
-1$. Conclusion~: $-1$ est un carré mod $p$ ssi $p = 2$ ou $p
\equiv 1 \pmod 4$.

\textbf{8.} Avec $x^2 \equiv -1$~: $p \mid x^2 + 1 = (x +
\iu)(x - \iu)$. Si $p$ était un premier de Gauss il diviserait
un facteur~; mais $\frac xp \pm \frac \iu p \notin \Z[\iu]$.
Donc $p$ n'est pas un premier de Gauss~; par la dichotomie
(question 4) --- $p$ non premier signifie la seconde branche ---
$p = a^2 + b^2$.

\textbf{9.} Les carrés sont $\equiv 0$ ou $1 \pmod 4$, donc $a^2
+ b^2 \in \{0, 1, 2\} \pmod 4$~: un premier $p \equiv 3 \pmod 4$
n'est pas une somme de deux carrés. Par la question 4, la
branche $N(\pi) = p$ ($p = a^2+b^2$) est impossible~: $p$ reste
un premier de Gauss.

\textbf{10.} $(1 + \iu)^2 = 2\iu$, donc $2 = -\iu(1 + \iu)^2$~;
et $N(1 + \iu) = 2$ est premier, donc $1 + \iu$ est un premier
de Gauss (question 2).

\textbf{11.} Tout premier de Gauss divise exactement un nombre
premier $p$ (question 3)~; listage par cas~: $p = 2$ donne les
\omterm{def:b3:rings:divisibility}{associés} de $1 + \iu$~; $p \equiv 3 \pmod 4$ donne $p$ lui-même
(question 9)~; $p \equiv 1 \pmod 4$ donne la paire $\pi,
\bar\pi$ de norme $p$ (questions 4 et 8). La paire est
authentique~: $\bar\pi \in \{\pm\pi, \pm\iu\pi\}$ forcerait, en
écrivant $\pi = a + \iu b$, soit $b = 0$, $a = 0$, soit $a =
\pm b$, donnant $p = a^2 + b^2 \in \{a^2, 2a^2\}$ --- impossible
pour un premier impair. Vérification~: $5 = (2 + \iu)(2 -
\iu)$, $N(2\pm\iu) = 5$~; $3$~: premier de norme $9$.

\textbf{12.} Écrire $n = 2^{\alpha}\prod_i p_i^{\beta_i}
\prod_j q_j^{2\gamma_j}$ avec $p_i \equiv 1$, $q_j \equiv 3
\pmod 4$. Chaque facteur est une somme de deux carrés~: $2 =
1^2 + 1^2$~; $p_i = a^2 + b^2$ (question 8)~; $q_j^{2\gamma_j}
= (q_j^{\gamma_j})^2 + 0^2$. L'identité de Brahmagupta
(question 1) propage la propriété au produit $n$.

\textbf{13.} Soit $n = a^2 + b^2 = N(a + \iu b)$ et $q \equiv 3
\pmod 4$, $q \mid n$. Le premier de Gauss $q$ (question 9)
divise $(a + \iu b)(a - \iu b)$, donc l'un des deux facteurs ---
disons $q \mid a + \iu b$ (l'autre cas est identique). Mais alors
$\frac{a + \iu b}{q} = \frac aq + \iu \frac bq \in \Z[\iu]$ se
lit comme $q \mid a$ \emph{et} $q \mid b$ dans $\Z$. D'où $q^2
\mid n$ et $\frac n{q^2} = \bigl(\frac aq\bigr)^2 + \bigl(\frac
bq\bigr)^2$. Par récurrence forte sur $n$, l'exposant de $q$ dans
$n/q^2$ est pair~; celui de $n$ l'est aussi.

\textbf{14.} \emph{Théorème (Fermat).} Un entier positif est une
somme de deux carrés si et seulement si tout premier $\equiv 3
\pmod 4$ y apparaît avec un exposant pair. --- $2025 = 3^4
\cdot 5^2$~: exposant de $3$ pair, oui ($2025 = 45^2 + 0^2 =
27^2 + 36^2$). $2026 = 2 \cdot 1013$ avec $1013 \equiv 1 \pmod
4$ premier~: oui ($1013 = 22^2 + 23^2$, et Brahmagupta avec $2 =
1^2+1^2$~: $2026 = (22 - 23)^2 + (22 + 23)^2 = 1^2 + 45^2$).
$2027$ est un premier $\equiv 3 \pmod 4$~: non.

\textbf{15.} Soit $p = a^2 + b^2 = c^2 + d^2$ avec entiers
positifs, $p \equiv 1 \pmod 4$, et $\pi$ un premier de Gauss
avec $p = \pi\bar\pi$ (question 4). $a + \iu b$ et $c + \iu d$
ont tous deux norme $p$, donc sont des premiers de Gauss
(question 2) divisant $p = (a+\iu b)(a - \iu b)$~; par unicité
de la factorisation, $c + \iu d$ est associé de $a + \iu b$ ou
de $a - \iu b$~:
\[
c + \iu d \in \{\pm(a \pm \iu b),\ \pm\iu(a \pm \iu b)\}
= \{\pm a \pm \iu b,\ \pm b \pm \iu a\}.
\]
La positivité de $c, d$ laisse $c + \iu d \in \{a + \iu b, b +
\iu a\}$~: $\{c, d\} = \{a, b\}$.

\textbf{16.} Si $n = a^2 - b^2 = (a-b)(a+b)$~: les deux facteurs
ont la même parité, donc $n$ est impair (tous deux impairs) ou
divisible par $4$ (tous deux pairs) --- jamais $\equiv 2 \pmod
4$. Réciproquement, $n$ impair~: $n = \bigl(\frac{n+1}2\bigr)^2
- \bigl(\frac{n-1}2\bigr)^2$~; $n = 4m$~: $n = (m+1)^2 -
(m-1)^2$. La réponse est une bare condition de congruence, avec
une identité d'une ligne derrière~: les différences de carrés
ne portent aucune profondeur arithmétique, et le contraste avec
les sommes est tout le point de ce problème.

\textbf{17.} $(a, b) \mapsto z = a + \iu b$ est une bijection
entre représentations et $\{z : N(z) = n\}$. Factoriser $z$ dans
l'\omterm{def:b3:rings:pidufd}{UFD} $\Z[\iu]$ en utilisant la classification (question 11)~: à
une unité près, $z = (1+\iu)^{a'}\prod_j\pi_j^{s_j}\bar\pi_j^{t_j}
\prod_kq_k^{u_k}$, et en prenant les normes ($N(1+\iu) = 2$,
$N(\pi_j) = N(\bar\pi_j) = p_j$, $N(q_k) = q_k^2$)~:
\[
n = 2^{a'}\prod_jp_j^{s_j + t_j}\prod_kq_k^{2u_k} .
\]
Appariement des exposants~: $a' = a$, $s_j + t_j = b_j$, $2u_k =
c_k$ --- solvable ssi tout $c_k$ est pair, et alors $u_k =
c_k/2$ est forcé tandis que $s_j \in \intint0{b_j}$ est libre.
Des données distinctes $(u, (s_j))$ donnent des $z$ non \omterm{def:b3:rings:divisibility}{associés}
de \emph{même} norme~; l'unité $u \in \{\pm1, \pm\iu\}$ ($4$
choix) énumère ensuite chaque classe d'\omterm{def:b3:rings:divisibility}{associés} sans répétition
(deux produits égaux violeraient l'unicité de la factorisation
--- $\pi_j$ et $\bar\pi_j$ sont non \omterm{def:b3:rings:divisibility}{associés} car $p_j =
\pi_j\bar\pi_j$ n'est pas ramifié). Total~: $r_2(n) =
4\prod_j(b_j + 1)$, et $0$ si un certain $c_k$ est impair.

\textbf{18.} $\chi(dd') = \chi(d)\chi(d')$ se vérifie mod $4$
(impair $\times$ impair couvre les quatre cas de signe~; tout
pair donne $0 = 0$). Pour $m, n$ premiers entre eux, les
diviseurs de $mn$ sont de façon unique $d = d_1d_2$ avec $d_1
\mid m$, $d_2 \mid n$~: $\sum_{d \mid mn}\chi(d) =
\bigl(\sum_{d_1\mid m}\chi(d_1)\bigr)\bigl(\sum_{d_2\mid
n}\chi(d_2)\bigr)$~: multiplicatif. Puissances de premiers~:
sur $2^a$, seul $d = 1$ est impair~: somme $= 1$. Sur $p^b$
avec $p \equiv 1$~: tous les $\chi(p^i) = 1$, somme $= b + 1$.
Sur $q^c$ avec $q \equiv 3$~: $\chi(q^i) = (-1)^i$, somme
alternée $= 1$ ($c$ pair) ou $0$ ($c$ impair).

\textbf{19.} Les deux fonctions multiplicatives $\frac14r_2$
(question 17) et $\sum_{d\mid n}\chi(d)$ (question 18)
s'accordent sur toutes les puissances de premiers --- $1$ sur
$2^a$~; $b + 1$ sur $p^b$~; $\mathbf 1_{c\ \mathrm{pair}}$ sur
$q^c$ --- donc s'accordent partout~: formule de Jacobi, avec
$\sum_{d\mid n}\chi(d) = d_1(n) - d_3(n)$ en triant les
diviseurs. Vérifications~: $r_2(3) = 0 = 4(1 - 1)$~; $r_2(5) =
8 = 4(2 - 0)$ ($(\pm1,\pm2), (\pm2,\pm1)$)~; $r_2(9) = 4 =
4(2 - 1)$ (diviseurs $1, 9 \equiv 1$~; $3 \equiv 3$~;
représentations $(\pm3, 0), (0, \pm3)$)~; $r_2(25) = 12 =
4(3 - 0)$. Pour $65 = 5\cdot13$~: $r_2 = 4\cdot2\cdot2 = 16$,
de $65 = 1 + 64 = 16 + 49$~: les seize paires $(\pm1, \pm8),
(\pm8, \pm1), (\pm4, \pm7), (\pm7, \pm4)$.

\textbf{20.} $\sum_{n \leq x}r_2(n)$ compte les paires $(a, b)$
avec $0 < a^2 + b^2 \leq x$, i.e.\ les points de réseau du
disque fermé $D_{\sqrt x}$ moins l'origine. Assigner à chaque
point de réseau $P$ le carré unité $P + \intco01^2$~: ces
carrés pavent le plan. Tout carré attaché à un point de
$D_{\sqrt x}$ vit dans $D_{\sqrt x + \sqrt2}$, et tout carré
rencontrant $D_{\sqrt x - \sqrt 2}$ est attaché à un point de
$D_{\sqrt x}$ (le carré a diamètre $\sqrt 2$)~: en comparant
les aires,
\[
\pi(\sqrt x - \sqrt2)^2 \leq \#\{\text{points de réseau dans }
D_{\sqrt x}\} \leq \pi(\sqrt x + \sqrt 2)^2,
\]
et les deux bornes sont $\pi x + O(\sqrt x)$. Soustraire
l'origine ne change rien à cette précision.

\textbf{21.} Par Jacobi (question 19) et en échangeant l'ordre
de sommation ($n = dm$)~:
\[
\frac14\sum_{n\leq x}r_2(n) = \sum_{n \leq x}\sum_{d \mid
n}\chi(d) = \sum_{d \leq x}\chi(d)\,\#\{m : dm \leq x\}
= \sum_{d\leq x}\chi(d)\Bigl\lfloor\frac xd\Bigr\rfloor,
\]
qui vaut $\frac{\pi x}4 + O(\sqrt x)$ par la question 20.
Enlever les parties entières~: $\lfloor x/d\rfloor = x/d +
O(1)$, mais sommer $O(1)$ sur $d \leq x$ est trop grossier~; à
la place utiliser que les sommes partielles de $\chi$ sont
bornées ($0, 1, 1, 0$ cycliquement), donc par sommation d'Abel
$\sum_{d\leq x}\chi(d)\{x/d\}$, dont les termes se groupent en
paires $d \equiv 1, 3$, est $O(\sqrt x)$ ---
alternativement et plus simplement~: scinder en $\sqrt x$. Pour
$d \leq \sqrt x$, remplacer $\lfloor x/d\rfloor$ par $x/d +
O(1)$~: erreur $O(\sqrt x)$. Pour $d > \sqrt x$, $\lfloor
x/d\rfloor$ prend chaque valeur $v < \sqrt x$ sur un intervalle
de $d$ consécutifs, sur lequel la somme en $\chi$ est $O(1)$~:
erreur totale $O(\sqrt x)$ en sommant sur les $\leq \sqrt x$
valeurs de $v$, tandis que $\sum_{d > \sqrt x}\chi(d)\frac xd =
O(\sqrt x)$ par les queues de séries alternées
($x\sum_{d>\sqrt x}\chi(d)/d = x\,O(1/\sqrt x)$). D'où
\[
x\sum_{d \leq x}\frac{\chi(d)}d = \frac{\pi x}4 + O(\sqrt x),
\qquad\text{i.e.}\qquad
\sum_{d\leq x}\frac{\chi(d)}d = \frac\pi4 +
O\Bigl(\frac1{\sqrt x}\Bigr),
\]
et en laissant $x \to \infty$~: $1 - \frac13 + \frac15 - \dots
= \frac\pi4$.

\textbf{22.} Si $n \equiv 3 \pmod4$ était $a^2 + b^2$~: les
carrés sont $\equiv 0, 1 \pmod 4$, et $a^2 + b^2 \in \{0, 1,
2\}$ mod $4$ --- impossible. (Le critère de la question 17 dit
la même chose~: $n \equiv 3 \pmod 4$ force un certain premier
$\equiv 3$ à exposant impair.) Donc les entiers représentables
évitent une classe de résidus pleine~: densité $\leq
\frac34$. La moyenne $\pi$ de $r_2$ se concentre sur peu
d'entiers~: $n = \prod_{j\leq k}p_j$ (premiers distincts $\equiv
1 \bmod 4$) a $r_2(n) = 4\cdot2^k$ représentations --- un nombre
non borné --- donc un ensemble creux de $n$ peut porter toute la
moyenne, exactement comme le gain moyen d'une loterie coexiste
avec une perte presque sûre. Le résultat de Landau
$\#\{n \leq x \text{ représentable}\} \sim Cx/\sqrt{\log x}$
le confirme~: densité $0$, moyenne $\pi$.

\textbf{23.} \emph{Nécessité.} Soit $n = a^2 + b^2$ avec
$\gcd(a, b) = 1$. Si un premier $q \equiv 3 \pmod 4$ divisait
$n$, la question 13 montre $q \mid a$ et $q \mid b$~:
contradiction. Si $4 \mid n$~: les carrés sont $\equiv 0, 1
\pmod 4$, donc $a^2 + b^2 \equiv 0 \pmod 4$ force $a^2 \equiv
b^2 \equiv 0$, i.e.\ $a, b$ tous deux pairs~: contradiction.
\emph{Suffisance.} Écrire $n = 2^{\alpha}\prod_jp_j^{b_j}$ avec
$\alpha \leq 1$ et $p_j \equiv 1 \pmod 4$, et poser $z =
(1+\iu)^{\alpha}\prod_j \pi_j^{b_j} = a + \iu b$, de norme $n$.
Supposer qu'un premier $t$ divise $\gcd(a, b)$~; alors $t \mid z$
dans $\Z[\iu]$. Si $t \equiv 3 \pmod 4$~: $t \mid N(z) = n$,
exclu. Si $t \equiv 1 \pmod 4$~: $t = \pi_t\bar\pi_t$, donc
$\bar\pi_t \mid z$~; mais la factorisation de $z$ ne contient
pas de premier conjugué ($\pi_j$ et $\bar\pi_j$ sont non
\omterm{def:b3:rings:divisibility}{associés}, question 17), contredisant l'unicité de la
factorisation. Si $t = 2 = -\iu(1+\iu)^2$~: alors $(1+\iu)^2
\mid z$, forçant $\alpha \geq 2$, exclu. D'où $\gcd(a, b) =
1$~: la représentation est primitive.

\textbf{24.} $a$ est impair ($\gcd(a, b) = 1$, $b$ pair), donc
$c^2 = a^2 + b^2$ est impair et $c$ est impair. Soit $\delta$
un diviseur premier de Gauss commun de $a + \iu b$ et $a - \iu
b$~: il divise leur somme $2a$ et leur différence $2\iu b$, donc
$2a$ et $2b$~; une relation de Bézout $ua + vb = 1$ donne alors
$\delta \mid 2$, donc $\delta$ est associé à $1 + \iu$ et
$N(\delta) = 2$ divise $N(a + \iu b) = c^2$, qui est impair~:
contradiction. Donc $a + \iu b$ et $a - \iu b$ sont premiers
entre eux de produit $c^2$~; dans l'\omterm{def:b3:rings:pidufd}{UFD} $\Z[\iu]$, chaque
premier de Gauss de $c^2$ apparaît à exposant pair et se scinde
entièrement dans l'un des deux facteurs premiers entre eux,
d'où $a + \iu b = u(m + \iu n)^2 = u\bigl(m^2 - n^2 + 2\iu
mn\bigr)$ avec $u$ une unité. Les choix $u = \pm\iu$ rendent la
partie réelle $\mp 2mn$ paire --- impossible, $a$ est impair.
Les choix $u = \pm1$ donnent, après ajustement des signes de
$m, n$ et échange de leurs noms pour tout rendre positif, $a =
m^2 - n^2$, $b = 2mn$ avec $m > n \geq 1$~; et $c^2 = N(m + \iu
n)^2$ donne $c = m^2 + n^2$. Un diviseur commun de $m$ et $n$
diviserait $a$ et $b$~: $\gcd(m, n) = 1$~; et $m \equiv n \pmod
2$ rendrait $a$ pair~: parités opposées. Vérifications~: $(m,
n) = (2, 1)$ donne $(3, 4, 5)$~; $(m, n) = (5, 2)$ donne $(25 -
4, 20, 25 + 4) = (21, 20, 29)$, et $441 + 400 = 841 = 29^2$.

\textbf{25.} La formule de Jacobi $r_2(n) = 4(d_1(n) - d_3(n))$
donne, pour $n = 1, \dots, 25$~:
\[
4,\ 4,\ 0,\ 4,\ 8,\ 0,\ 0,\ 4,\ 4,\ 8,\ 0,\ 0,\ 8,\ 0,\ 0,\
4,\ 8,\ 4,\ 0,\ 8,\ 0,\ 0,\ 0,\ 0,\ 12,
\]
non nuls exactement en $n = 1, 2, 4, 5, 8, 9, 10, 13, 16, 17,
18, 20, 25$ (par exemple $r_2(15) = 0$~: les diviseurs $1, 5
\equiv 1$ et $3, 15 \equiv 3$ s'équilibrent~; $r_2(20) = 8$~:
diviseurs $1, 5 \equiv 1$, aucun $\equiv 3$). Le total est $4 +
4 + 4 + 8 + 4 + 4 + 8 + 8 + 4 + 8 + 4 + 8 + 12 = 80$. Le côté
diviseurs~: les $d$ impairs $\leq 25$ contribuent
\[
25 - 8 + 5 - 3 + 2 - 2 + 1 - 1 + 1 - 1 + 1 - 1 + 1 = 20,
\]
en lisant $\chi(d)\lfloor 25/d\rfloor$ pour $d = 1, 3, 5,
\dots, 25$~; et $4 \cdot 20 = 80$, comme prédit par l'identité
de la question 21. Points de réseau du disque fermé de rayon
$5$~: les $80$ points avec $1 \leq a^2 + b^2 \leq 25$ plus
l'origine, i.e.\ $81$~; et $\pi x = 25\pi \approx 78{,}54$, une
erreur d'environ $2{,}46$, confortablement dans la bande
$O(\sqrt x)$ de la question 20 ($\sqrt x = 5$).
\end{solution}
